\documentclass[journal]{IEEEtran}

\usepackage[T1]{fontenc}
\usepackage{lmodern}
\usepackage{amsmath,amssymb}
\usepackage{booktabs}
\usepackage{graphicx}
\usepackage[hidelinks]{hyperref}
\usepackage{url}

\newcommand{\systemname}{TPG-Fact}
\newcommand{\datasetname}{US-TPMFC}

\title{Temporal and Evidence-Source Provenance-Aware Multimodal Misinformation Detection with Calibrated Abstention}

\author{Sabbir Hossain\\
Department of Computer Science\\
American International University-Bangladesh\\
Dhaka, Bangladesh\\
\href{mailto:26-94030-1@student.aiub.edu}{26-94030-1@student.aiub.edu}}

\begin{document}
\maketitle

\begin{abstract}
Web evidence can mislead a multimodal fact-checker when it was published too late, copied across several sites, or taken from sources that depend on one another. \systemname{} addresses these problems by enforcing a decision-time cutoff and grouping syndicated documents before aggregation. It represents claims, documents, images, publishers, links, and publication order in a heterogeneous graph. A hierarchical classifier first estimates whether the retrieved material is sufficient. When it is, a second head predicts whether the claim is supported or refuted; otherwise, the system returns insufficient evidence. Temperature scaling and a validation-selected threshold support selective abstention. The accompanying \datasetname{} benchmark uses United States ClaimReview cases and evidence recovered from open web archives. Training covers 2020--2023, validation uses 2024, and testing uses 2025--2026, with related events, near-duplicate claims, and reused images kept within one partition. The evaluation compares the method with matched multimodal retrieval baselines and tests temporal leakage, syndicated evidence, calibration, and selective risk.
\end{abstract}

\begin{IEEEkeywords}
multimodal misinformation, automated fact-checking, evidence provenance, temporal retrieval, selective prediction, calibrated abstention
\end{IEEEkeywords}

\section{Introduction}

An authentic photograph can still mislead when it is paired with the wrong caption, date, place, or event. Models that focus on manipulation artifacts or surface-level image--text compatibility may miss this kind of false context. External evidence can help, but it creates its own problems. Some documents appear only after the decision deadline, search results may repeat one syndicated report, and a model can sound certain even when the available evidence is weak.

Several multimodal systems address parts of this problem. KAMP learns from text, images, and background knowledge graphs \cite{zhang2025kamp}. RED-DOT reranks retrieved multimodal evidence by relevance \cite{papadopoulos2025reddot}, while COVE recovers an image's original context before judging its caption \cite{tonglet2025cove}. Other systems use multimodal retrieval-augmented generation to reason over web sources \cite{khaliq2024ragar,kangur2025multireflect}. None of these operations alone establishes that a document existed at decision time or that a list of results represents independent sources.

We ask whether time-valid evidence from independent sources improves future-event misinformation detection and reduces overconfident decisions relative to ordinary multimodal RAG. Each case is an English image--text claim about a United States event. The system returns supported, refuted, or insufficient evidence, together with the evidence identifiers and cutoff used for the decision.

The study design has four parts. \datasetname{} provides a temporal benchmark built from professional fact-check metadata and open archives. A provenance graph records publication order, duplication, hyperlinks, publishers, entities, events, and shared media. The detector evaluates evidence sufficiency before it predicts veracity. The evaluation then measures future-event performance, temporal leakage, sensitivity to copied sources, calibration, and selective risk.


\section{Related Work}

\subsection{Multimodal Misinformation Detection}

NewsCLIPpings introduced a large benchmark of authentic news images paired with deliberately mismatched captions \cite{luo2021newsclippings}. VERITE later addressed unimodal bias and topic imbalance in real-world evaluation \cite{papadopoulos2024verite}. KAMP uses single-modality, cross-modal, and knowledge-content objectives during multimodal pretraining \cite{zhang2025kamp}. Other work uses synthetic examples to improve transfer to real misinformation cases \cite{zeng2024synthetic}. Event-Radar represents visual manipulation, textual emotion, and event-level inconsistency in an event graph \cite{ma2024eventradar}. These methods provide the main no-retrieval and event-aware baselines for our study.

\subsection{Evidence-Grounded Fact-Checking}

MOCHEG treats multimodal fact-checking as a sequence of evidence retrieval, claim verification, and explanation generation \cite{yao2023mocheg}. MetaSumPerceiver summarizes evidence drawn from several multimodal documents \cite{chen2024metasum}. RED-DOT accounts for irrelevant web evidence through relevance-aware reranking \cite{papadopoulos2025reddot}. RAGAR and MultiReflect combine multimodal language models with retrieval and iterative reasoning \cite{khaliq2024ragar,kangur2025multireflect}, and SNIFFER retrieves external material for out-of-context detection \cite{qi2024sniffer}. Our method builds on this retrieval setting while enforcing publication-time cutoffs and aggregating evidence by source dependency.

\subsection{Image Context, Time, and Evidence Dependence}

The 5Pils system recovers an image's original context from web evidence \cite{tonglet2024image}, and COVE predicts that context before assessing a caption \cite{tonglet2025cove}. Chronological distribution shift can reduce the accuracy of multimodal misinformation models \cite{stepanova2023temporal}. ChronoFact handles complex temporal claims with structured timeline reasoning \cite{barik2025chronofact}. Source information also affects how RAG-based fact verification resolves conflicting evidence \cite{ge2025conflicting}. We represent publication time and source dependence in the same evidence graph.


\section{Task and Benchmark Construction}

\subsection{Task Definition}

Each instance contains claim text $x_t$, an image $x_v$, a decision cutoff $\tau$, and an evidence collection $E_{\leq\tau}$. The output is $y \in \{\text{supported},\text{refuted},\text{insufficient}\}$. Documents published after $\tau$ are unavailable to the model. The event is the independent unit for data splitting and uncertainty analysis.

\subsection{Candidate Acquisition}

Corpus construction starts with English ClaimReview records about United States events from 2020--2026. GDELT provides discovery metadata, and Common Crawl provides archived page observations. We retain cases with an identifiable claim time, a recoverable image or archive, an unambiguous professional rating, and enough timestamp information to enforce the cutoff. The retrieval pool excludes fact-check domains and documents that contain the final verdict.

\subsection{Label Policy}

We normalize only explicit site ratings. True or correct maps to supported. False, fake, or miscaptioned maps to refuted. Unproven, unsubstantiated, or no evidence maps to insufficient evidence. We exclude mixed, partially true, satire, opinion, and ambiguous ratings. The labels come from professional fact-checkers; automatically extracted evidence links are treated as silver evidence.

\subsection{Temporal Partitioning and Release}

Training covers 2020--2023, validation uses 2024, and final testing uses 2025--2026. Events, near-duplicate claims, and images with matching perceptual hashes remain in one partition. The release package contains URLs, cryptographic and perceptual hashes, timestamps, normalized labels, graph annotations, retrieval queries, and split manifests. It does not redistribute copyrighted articles or images.

We treat \datasetname{} as a benchmark only if the filtered corpus contains at least 2,000 cases, 300 cases per label, 50 test events, and two independent evidence clusters for at least 70\% of retained claims. If the collection misses any threshold, we use it as a temporal and provenance stress test instead.


\section{Temporal-Provenance Method}

\subsection{Time-Constrained Multimodal Retrieval}

For each claim, retrieval begins with named entities, temporal expressions, OCR text, and an image caption. BM25 and BGE-M3 retrieve text candidates, while SigLIP2 embeddings and perceptual hashes retrieve images. Reciprocal-rank fusion merges the lexical, dense, and visual rankings. A reranker processes the top 100 candidates, and the system keeps the 20 highest-ranked items that satisfy the publication-time cutoff.

\subsection{Evidence-Source Provenance Graph}

The graph contains nodes for claims, images, documents, publishers, entities, and events. Typed edges record retrieval relevance, image containment, near-duplicate content, hyperlinks, publisher membership, entity or event mentions, and temporal order. We estimate document clusters from exact hashes, MinHash-style lexical similarity, semantic similarity, perceptual image distance, hyperlinks, and publication order. The earliest observed document is an origin candidate because the available records may not reveal the original author.

Let $C(E)$ denote the partition of evidence documents into dependency clusters. The model first aggregates documents within a cluster and then combines information across clusters. As a result, repeated copies of one report do not receive the same weight as several independent sources.

\subsection{Hierarchical Decision and Abstention}

A heterogeneous graph transformer encodes the temporal-provenance graph for the claim. The sufficiency head estimates $p(s=1\mid x,E_{\leq\tau})$. Sufficient evidence passes to a veracity head that predicts supported or refuted; the remaining cases receive the insufficient-evidence label. We fit temperature scaling to validation predictions and choose an abstention threshold on the same split to target 80\% answer coverage. Each prediction records class probabilities, evidence identifiers, source clusters, and the enforced cutoff.


\section{Experimental Protocol}

\subsection{Baselines}

The controlled comparison includes majority prediction, text-only classification, image--text fusion without retrieval, and an open multimodal language model under zero- and few-shot prompting. Retrieval baselines cover vanilla multimodal RAG, time-filtered RAG without provenance, provenance-aware RAG without time filtering, and RED-DOT where its implementation is compatible with the protocol. Every method uses the same splits, candidate corpus, evaluation code, and 20-trial validation budget.

\subsection{Metrics and Statistical Analysis}

The primary metric is macro-F1 on the 2025--2026 event-disjoint test set. Secondary metrics are per-class F1, balanced accuracy, macro one-versus-rest AUROC, Brier score, negative log-likelihood, expected calibration error, area under the risk--coverage curve, risk at 80\% coverage, Recall@10, nDCG@10, latency, GPU memory, and retrieved-document count. Final configurations are evaluated with five seeds. Differences use 10,000 paired event-level bootstrap replicates, with Holm correction across the primary comparison family.

\subsection{Ablation and Stress Tests}

The ablations remove temporal filtering, dependency clustering, temporal edge features, graph aggregation, text evidence, visual evidence, the sufficiency gate, and evidence-dropout training one at a time. The stress tests add syndicated copies and post-cutoff documents that reveal the verdict. They also remove the earliest observed source, introduce conflicts between dependent and independent evidence, swap captions, and reuse historical images.

\subsection{Selection and Falsification Rules}

We select hyperparameters and the abstention threshold using validation data only. Evidence-source provenance must improve either future-event performance or resistance to duplicate injection to support a performance claim. Calibrated abstention must reduce selective risk at matched coverage. Gains that depend on post-cutoff documents count as leakage.


\section{Evaluation Questions}

The first analysis tests whether excluding post-cutoff documents changes retrieval quality and future-event classification. The full method is compared with vanilla multimodal RAG and with provenance-aware RAG that does not enforce a time cutoff. Event-level paired bootstrap intervals quantify the differences.

The second analysis measures how source dependence affects a verdict. It compares cluster-level aggregation with document-level aggregation before and after injecting one, five, and ten syndicated copies of the same source. Changes in macro-F1, confidence, Brier score, and selective risk show whether repeated evidence is being counted as independent corroboration.

The third analysis examines the sufficiency gate and calibrated abstention. A flat three-class classifier and an uncalibrated confidence threshold provide matched controls. The comparison uses insufficient-evidence F1, area under the risk--coverage curve, and risk at 80\% coverage.


\section{Discussion, Ethics, and Limitations}

The main methodological risk is mistaking repeated web coverage for independent corroboration. Dependency clusters can only use observable content, links, domains, and timestamps. They cannot prove that two publishers worked independently or identify the original author with certainty. Archive timestamps show when a page was observed, which may be later than its first publication.

The corpus inherits professional fact-check labels and the editorial policies behind them. Restricting the mapping to explicit ratings reduces ambiguity, but it can also create selection bias and class imbalance. The focus on United States events and English sources limits geographic and linguistic generalization. Common Crawl and GDELT do not cover the whole web, and blocked or deleted pages may differ systematically from the pages that remain available.

The release will not redistribute copyrighted articles or images without permission. It will contain URLs, hashes, derived metadata, split manifests, and collection code. The system is meant for research and assisted review, not autonomous censorship or content removal. An insufficient-evidence prediction calls for further investigation; it does not mean that the claim is false.


\section{Conclusion}

This study examines whether a multimodal fact-checker benefits from knowing when evidence appeared and when several documents came from the same source. \systemname{} enforces a decision-time cutoff, groups dependent evidence before aggregation, checks whether the retrieved material is sufficient, and abstains when confidence is too low. The evaluation separates ordinary future-event performance from gains caused by temporal leakage and repeated documents. Its ablations also test whether the graph, source clusters, and sufficiency gate earn their place in the model. The final claims depend on controlled comparisons; a large stack of web documents does not count as corroboration on its own.



\bibliographystyle{IEEEtran}
\bibliography{references}

\end{document}

